\chapter{Fuerzas y movimiento}\label{ch:g11:forces-and-motion}

Lanza una llave inglesa girando de un lado a otro de la habitación: los
extremos se bambolean sin orden, y sin embargo un punto de su interior
describe una parábola perfecta. El \cref{ch:g10:inertia} resolvió qué
ocurre cuando las \omterm{def:g10:inertia:force}{fuerzas} \omterm{def:g10:inertia:compensate}{se compensan} --- nada. Este capítulo pregunta
qué ocurre cuando no lo hacen: la \omterm{def:g10:inertia:force}{fuerza} sobrante \emph{cambia la
velocidad} --- acelera los cuerpos, los frena, los hace girar --- y dice
cuánto, en proporciones si todavía no en fórmulas.

\section{El centro de masas}

\begin{definition}[Centro de masas]\label{def:g11:forces-and-motion:com}
El \emph{centro de masas}\index{centro de masas} de un cuerpo o de un
sistema es su punto de equilibrio --- la posición media de su masa: el
centro geométrico de un cuerpo homogéneo y simétrico; para dos masas
puntuales, el punto del segmento que las une situado a distancias en
razón \emph{inversa} a las masas,
$\frac{d_1}{d_2} = \frac{m_2}{m_1}$ --- más cerca de la más pesada.
\end{definition}

\begin{example}[Una pesa descompensada]\label{ex:g11:forces-and-motion:dumbbell}
Hay bolas de \qty{2.0}{kg} y \qty{1.0}{kg} en los extremos de una barra
ligera de \qty{0.90}{m}. El \omterm{def:g11:forces-and-motion:com}{centro de masas} la divide en la razón
inversa $d_1/d_2 = 1.0/2.0$: a $\qty{0.30}{m}$ de la bola pesada y a
$\qty{0.60}{m}$ de la ligera --- equilibra la barra ahí.
\end{example}

\begin{remark}[La llave lanzada]\label{rem:g11:forces-and-motion:wrench}
Graba una llave inglesa dando vueltas: todos sus puntos siguen una curva
complicada llena de bucles --- todos salvo el \omterm{def:g11:forces-and-motion:com}{centro de masas}, que
describe la parábola limpia de una pelota lanzada. El movimiento del
\omterm{def:g11:forces-and-motion:com}{centro de masas} es sencillo; el resto es giro \emph{a su alrededor} ---
y por eso la mecánica puede tratar un coche, un planeta o una gimnasta
como un solo punto.
\end{remark}

\begin{omfigure}
\begin{tikzpicture}[scale=1.0, font=\small]
  \draw[gray!60] (-0.4,0) -- (8.6,0);
  \draw[dashed, omDef, thick]
    plot[domain=0:8.2, samples=60] (\x, {0.9*\x - 0.11*\x*\x});
  \draw[very thick, omExo] ($(0.5,0.42)+(20:0.55)$) -- ($(0.5,0.42)-(20:0.55)$);
  \draw[very thick, omExo] ($(2.3,1.49)+(95:0.55)$) -- ($(2.3,1.49)-(95:0.55)$);
  \draw[very thick, omExo] ($(4.1,1.84)+(170:0.55)$) -- ($(4.1,1.84)-(170:0.55)$);
  \draw[very thick, omExo] ($(5.9,1.48)+(245:0.55)$) -- ($(5.9,1.48)-(245:0.55)$);
  \draw[very thick, omExo] ($(7.7,0.41)+(320:0.55)$) -- ($(7.7,0.41)-(320:0.55)$);
  \fill[omDef] (0.5,0.42) circle (2pt) (2.3,1.49) circle (2pt)
    (4.1,1.84) circle (2pt) (5.9,1.48) circle (2pt) (7.7,0.41) circle (2pt);
  \node[omDef] at (4.1,2.4) {centro de masas};
\end{tikzpicture}

{\small Una llave inglesa lanzada, a intervalos de tiempo iguales: la
llave gira, pero su \omterm{def:g11:forces-and-motion:com}{centro de masas} (los puntos) cabalga sobre una
parábola sencilla.}
\end{omfigure}

\section{Una fuerza neta cambia la velocidad}

\begin{definition}[Fuerza neta]\label{def:g11:forces-and-motion:netforce}
La \emph{fuerza neta}\index{fuerza neta} sobre un cuerpo es la suma
vectorial de todas las \omterm{def:g10:inertia:force}{fuerzas} que actúan sobre él,
$\vect F = \vect F_1 + \vect F_2 + \dots$ (punta con cola, como los
vectores del volumen de matemáticas); las \omterm{def:g10:inertia:force}{fuerzas} \omterm{def:g10:inertia:compensate}{se compensan}
exactamente cuando $\vect F = \vect 0$.
\end{definition}

\begin{theorem}[Efecto de una fuerza neta --- versión semicuantitativa]\label{thm:g11:forces-and-motion:secondlaw}
\index{segunda ley del movimiento}
Una \omterm{def:g11:forces-and-motion:netforce}{fuerza neta} $\vect F$ que actúa sobre un cuerpo de masa $m$ durante
un intervalo corto $\Delta t$ cambia la velocidad de su \omterm{def:g11:forces-and-motion:com}{centro de masas}
en un vector $\Delta \vect v$ que tiene la \emph{dirección y el sentido}
de $\vect F$, y cuyo módulo crece en proporción a $F$ y a $\Delta t$ y
disminuye en proporción a la masa:
\[
\Delta v \propto \frac{F \, \Delta t}{m} .
\]
\end{theorem}
\admitted

\begin{remark}
Duplica la \omterm{def:g10:inertia:force}{fuerza} o su duración: el doble de cambio; duplica la masa: la
mitad. El curso que viene afina esto hasta convertirlo en una ecuación
exacta, y el volumen del primer año pregunta en qué sistemas se cumple;
este año bastan las proporciones --- que ya gobiernan airbags, frenos y
\omterm{def:g10:universal-gravitation:satellite}{órbitas}.
\end{remark}

\begin{proposition}[Los tres casos canónicos]\label{prop:g11:forces-and-motion:cases}
Sea un cuerpo que se mueve con velocidad $\vect v$ bajo una \omterm{def:g11:forces-and-motion:netforce}{fuerza neta}
$\vect F$.
\begin{enumerate}
  \item $\vect F$ en el sentido de $\vect v$: \emph{acelera} en línea
        recta;
  \item $\vect F$ en sentido opuesto a $\vect v$: \emph{frena} en línea
        recta;
  \item $\vect F$ perpendicular a $\vect v$: \emph{gira} sin cambiar de
        rapidez.
\end{enumerate}
\end{proposition}

\begin{proof}
Se suma $\Delta\vect v$, dirigida según $\vect F$
(\cref{thm:g11:forces-and-motion:secondlaw}), punta con cola sobre
$\vect v$: en el sentido de $\vect v$ las flechas se apilan (más larga);
en sentido opuesto se cancelan en parte (más corta); perpendicular, cada
pequeño $\Delta \vect v$ inclina la flecha sin alargarla --- solo gira
la dirección.
\end{proof}

\begin{omfigure}
\begin{tikzpicture}[scale=0.9, font=\small]
  \draw[-Stealth, omDef, thick] (0,1.5) -- (1.4,1.5) node[above] {$\vect v$};
  \draw[-Stealth, omThm, very thick] (0,0.75) -- (0.8,0.75) node[above] {$\vect F$};
  \draw[-Stealth, omDef, thick] (0,0) -- (2.2,0) node[above] {$\vect v'$};
  \node at (1.1,-0.7) {acelera};
\end{tikzpicture}
\qquad
\begin{tikzpicture}[scale=0.9, font=\small]
  \draw[-Stealth, omDef, thick] (0,1.5) -- (2.2,1.5) node[above] {$\vect v$};
  \draw[-Stealth, omThm, very thick] (0.8,0.75) -- (0,0.75) node[above] {$\vect F$};
  \draw[-Stealth, omDef, thick] (0,0) -- (1.4,0) node[above] {$\vect v'$};
  \node at (1.1,-0.7) {frena};
\end{tikzpicture}
\qquad
\begin{tikzpicture}[scale=0.9, font=\small]
  \draw[-Stealth, omDef, thick] (0,1.5) -- (2.0,1.5) node[above] {$\vect v$};
  \draw[-Stealth, omThm, very thick] (1.0,0.4) -- (1.0,1.2) node[right] {$\vect F$};
  \draw[-Stealth, omDef, thick] (0,-0.5) -- (30:2.0)
    node[above] {$\vect v'$};
  \node at (1.1,-1.1) {gira};
\end{tikzpicture}

{\small Los tres casos canónicos: la \omterm{def:g11:forces-and-motion:netforce}{fuerza neta} arrastra la punta del
vector velocidad hacia su propio lado --- más larga, más corta o
girada.}
\end{omfigure}

\begin{example}[Un golpe lateral]\label{ex:g11:forces-and-motion:kick}
Un disco se desliza hacia el este a \qty{6.0}{m/s}; un golpe lateral
breve le añade $\Delta\vect v = \qty{8.0}{m/s}$ hacia el norte. La nueva
velocidad es la suma vectorial: de módulo
$\sqrt{6.0^2 + 8.0^2} = \qty{10.0}{m/s}$ y dirección
$\tan\theta = 8.0/6.0$, es decir $\theta \approx 53^\circ$ del este
hacia el norte --- la velocidad antigua no se borra: el cambio se le
suma.
\end{example}

\section{Comparar cambios: fuerza, tiempo y masa}

\begin{method}[Razonar con proporciones]\label{met:g11:forces-and-motion:ratio}
Para comparar dos situaciones, nunca hay que resolver una ecuación:
basta formar cocientes. El cociente de los cambios de velocidad es el
cociente de las $F$, por el cociente de los $\Delta t$, dividido entre
el cociente de las $m$. Después se lee qué factor ha cambiado y en
cuánto.
\end{method}

\begin{example}[El carrito cargado]\label{ex:g11:forces-and-motion:cart}
El mismo empujón, aplicado durante el mismo segundo, lanza un carrito de
la compra vacío de \qty{20}{kg} a \qty{1.2}{m/s}. Cargado hasta
\qty{60}{kg} --- el triple de masa, el mismo $F\Delta t$ --- sale a la
tercera parte, \qty{0.40}{m/s}. La masa es la terquedad de la materia.
\end{example}

\begin{remark}[Choques, airbags y rodillas flexionadas]\label{rem:g11:forces-and-motion:airbag}
En un choque, $\Delta v$ y $m$ no se negocian --- la velocidad del
ocupante tiene que llegar a cero ---, así que $F\,\Delta t$ está fijado.
La única libertad que queda es el \emph{tiempo}: alarga $\Delta t$ y la
\omterm{def:g10:inertia:force}{fuerza} se encoge en la misma proporción. Una cabeza que se detiene sobre
un airbag en \qty{0.08}{s} en lugar de sobre el volante en
\qty{0.01}{s} nota una \omterm{def:g10:inertia:force}{fuerza} media ocho veces menor; las zonas de
deformación, la espuma de un casco y las rodillas flexionadas hacen el
mismo truco.
\end{remark}

\section{El rozamiento: la fuerza que se opone al deslizamiento}

\begin{definition}[Rozamiento estático y cinético]\label{def:g11:forces-and-motion:friction}
El \omterm{def:g10:inertia:inventory}{rozamiento} entre superficies en contacto se opone a su
\emph{deslizamiento relativo}. El \emph{rozamiento
estático}\index{rozamiento!estático} actúa mientras no deslizan y se
ajusta solo, hasta un máximo, para cancelar lo que intente iniciar el
deslizamiento; el \emph{rozamiento
cinético}\index{rozamiento!cinético}, una vez empezado el
deslizamiento, es aproximadamente constante, se dirige en contra del
deslizamiento y suele ser \emph{más débil} que el máximo estático.
\end{definition}

\begin{example}[Empujar el armario]\label{ex:g11:forces-and-motion:wardrobe}
Empuja un armario con suavidad: no se mueve --- el \omterm{def:g11:forces-and-motion:friction}{rozamiento estático}
cancela tus \qty{150}{N}. Empuja más fuerte: hacia los \qty{400}{N} se
pone en marcha de un tirón y después desliza con menos esfuerzo --- ha
tomado el relevo el \omterm{def:g11:forces-and-motion:friction}{rozamiento cinético}, más débil.
\end{example}

\begin{remark}[Por qué funciona el ABS]\label{rem:g11:forces-and-motion:abs}
El punto de contacto de una rueda que gira no desliza sobre la
carretera: el neumático agarra con \omterm{def:g10:inertia:inventory}{rozamiento} \emph{estático}. Una rueda
bloqueada derrapa y solo obtiene el \omterm{def:g10:inertia:inventory}{rozamiento} \emph{cinético}, más
débil, que además se opone al derrape hagan lo que hagan las ruedas ---
y la dirección muere. El ABS suelta el freno en cuanto una rueda se
bloquea: más agarre, frenada más corta y un conductor que puede girar.
\end{remark}

\begin{omfigure}
\begin{tikzpicture}
\begin{axis}[omaxis, width=8.6cm, height=5.2cm, ybar, bar width=16pt,
  symbolic x coords={carretera seca, carretera mojada, nieve pisada}, xtick=data,
  xlabel={}, ylabel={distancia de frenado (\unit{m})}, ymin=0, ymax=185,
  nodes near coords,
  every node near coord/.append style={font=\small},
  enlarge x limits=0.25]
\addplot[draw=omDef, fill=omDef!25] coordinates
  {(carretera seca,40) (carretera mojada,70) (nieve pisada,160)};
\end{axis}
\end{tikzpicture}

{\small Distancias de frenado típicas desde \qty{90}{km/h}: el
\omterm{def:g10:inertia:inventory}{rozamiento} entre neumático y calzada es la única \omterm{def:g10:inertia:force}{fuerza} que frena el
coche, y el agua o la nieve lo recortan gravemente.}
\end{omfigure}

\section{Girar: el movimiento circular necesita una fuerza hacia dentro}

\begin{definition}[Movimiento circular uniforme]\label{def:g11:forces-and-motion:ucm}
Un cuerpo tiene un \emph{movimiento circular
uniforme}\index{movimiento circular uniforme} cuando su \omterm{def:g10:relative-motion:trajectory}{trayectoria} es
una circunferencia recorrida con rapidez constante. La velocidad,
tangente a la circunferencia, cambia de dirección en cada instante ---
de modo que nunca es constante.
\end{definition}

\begin{proposition}[La fuerza hacia dentro]\label{prop:g11:forces-and-motion:inward}
Un cuerpo con \omterm{def:g11:forces-and-motion:ucm}{movimiento circular uniforme} está sometido a una \omterm{def:g11:forces-and-motion:netforce}{fuerza
neta} que nunca es nula: en cada instante apunta desde el cuerpo hacia el
centro de la circunferencia.
\end{proposition}

\begin{proof}
La rapidez es constante, así que la \omterm{def:g11:forces-and-motion:netforce}{fuerza neta} no tiene componente
según $\vect v$ (los casos 1 y 2 de la
\cref{prop:g11:forces-and-motion:cases}): es perpendicular a la
tangente, o sea, radial --- y hacia dentro, porque cada
$\Delta \vect v$ curva el camino hacia el centro.
\end{proof}

\begin{omfigure}
\begin{tikzpicture}[scale=1.0, font=\small]
  \draw[gray!50, dashed] (0,0) circle (1.7);
  \fill (0,0) circle (1.5pt) node[below] {mano};
  \draw[gray!70] (0,0) -- (90:1.7) (0,0) -- (210:1.7) (0,0) -- (330:1.7);
  \fill[omExo] (90:1.7) circle (2.5pt);
  \fill[omExo] (210:1.7) circle (2.5pt);
  \fill[omExo] (330:1.7) circle (2.5pt);
  \draw[-Stealth, omThm, very thick] (90:1.7) -- (90:0.75) node[pos=0.35, right] {$\vect F$};
  \draw[-Stealth, omThm, very thick] (210:1.7) -- (210:0.75);
  \draw[-Stealth, omThm, very thick] (330:1.7) -- (330:0.75);
  \draw[-Stealth, omDef, thick] (90:1.7) -- ++(180:1.2) node[left] {$\vect v$};
  \draw[-Stealth, omDef, thick] (210:1.7) -- ++(300:1.2);
  \draw[-Stealth, omDef, thick] (330:1.7) -- ++(60:1.2);
\end{tikzpicture}

{\small Una bola volteada con una cuerda: la velocidad es tangente y el
tirón apunta al centro. Si se corta la cuerda, la bola sale por la
tangente.}
\end{omfigure}

\begin{example}[La cuerda, la curva y la Luna]\label{ex:g11:forces-and-motion:moon}
Todo giro estable esconde una \omterm{def:g10:inertia:force}{fuerza} hacia dentro. La bola volteada: la
tensión de la cuerda. El coche en una curva: el \omterm{def:g11:forces-and-motion:friction}{rozamiento estático}
lateral de la calzada sobre los neumáticos --- que desaparece sobre
hielo, y entonces el coche sigue recto. La Luna: el tirón gravitatorio
de la Tierra (\cref{ch:g10:universal-gravitation}), que la hace girar en
\num{27.3} días --- no la sostiene \emph{en alto}, sino que tira de ella
de lado hacia un giro perpetuo, cayendo alrededor de la Tierra; el curso
que viene convierte esto en la mecánica de los \omterm{def:g10:universal-gravitation:satellite}{satélites}.
\end{example}

\section{Ejercicios}

\begin{exercise}[$\star$]\label{exo:g11:forces-and-motion:1}
¿\omterm{def:g11:forces-and-motion:netforce}{Fuerza neta} nula o no nula? Justifícalo a partir de la velocidad: (a)
un coche que circula recto a \qty{90}{km/h} constantes; (b) que frena en
línea recta; (c) que toma una curva a \qty{50}{km/h} constantes; (d) un
paracaidista a velocidad límite.
\end{exercise}

\begin{exercise}[$\star$]\label{exo:g11:forces-and-motion:2}
Hay bolas de \qty{3.0}{kg} y \qty{1.0}{kg} en los extremos de una barra
ligera de \qty{0.80}{m}. Localiza el \omterm{def:g11:forces-and-motion:com}{centro de masas}; ¿dónde estaría si
las masas fueran iguales?
\end{exercise}

\begin{exercise}[$\star$]\label{exo:g11:forces-and-motion:3}
Se da el mismo empujón breve a un carro de equipaje vacío de
\qty{15}{kg} y a ese mismo carro cargado hasta \qty{45}{kg}. Compara los
cambios de velocidad; ¿qué empujón le daría al carro cargado el
$\Delta v$ del vacío?
\end{exercise}

\begin{exercise}[$\star$]\label{exo:g11:forces-and-motion:4}
Un tranvía que va a \qty{10}{m/s} frena; un segundo después va a
\qty{8}{m/s} en el mismo sentido. Da $\Delta \vect v$ y la dirección de
la \omterm{def:g11:forces-and-motion:netforce}{fuerza neta}. ¿Qué caso de la
\cref{prop:g11:forces-and-motion:cases} es este?
\end{exercise}

\begin{exercise}[$\star$]\label{exo:g11:forces-and-motion:5}
¿\omterm{def:g11:forces-and-motion:friction}{Rozamiento estático} o cinético --- y dirigido hacia dónde? (a) Una caja
que resiste tu empujón (demasiado suave); (b) esa misma caja una vez que
desliza; (c) un libro sobre una bandeja ligeramente inclinada que no
desliza.
\end{exercise}

\begin{exercise}[$\star\star$]\label{exo:g11:forces-and-motion:6}
Una \omterm{def:g11:forces-and-motion:netforce}{fuerza neta} $F$ aplicada durante $\Delta t$ a una masa $m$ produce
$\Delta v = \qty{4.0}{m/s}$. Predice $\Delta v$ para: (a)
$2F, \Delta t, m$; (b) $F, \Delta t/2, m$; (c) $F, \Delta t, 2m$; (d)
$3F, 2\Delta t, 6m$.
\end{exercise}

\begin{exercise}[$\star\star$]\label{exo:g11:forces-and-motion:7}
Un disco desliza hacia el este a \qty{8.0}{m/s}; un golpe le añade
$\Delta \vect v = \qty{6.0}{m/s}$ hacia el norte. Calcula la nueva
rapidez y la dirección. ¿Por qué la rapidez no es
$8.0 + 6.0 = \qty{14}{m/s}$?
\end{exercise}

\begin{exercise}[$\star\star$]\label{exo:g11:forces-and-motion:8}
Una portera para el mismo disparo de dos maneras: con los brazos rígidos
(el balón se detiene en \qty{0.02}{s}) o ``acompañando'' el balón
(\qty{0.08}{s}). Compara las \omterm{def:g10:inertia:force}{fuerzas} medias y di qué magnitudes estaban
fijadas antes de que ella eligiera.
\end{exercise}

\begin{exercise}[$\star\star$]\label{exo:g11:forces-and-motion:9}
¿Qué \omterm{def:g10:inertia:force}{fuerza} exterior impulsa hacia adelante a un velocista en la salida?
¿Por qué ayudan las zapatillas de clavos --- y por qué nadie puede
acelerar, ni siquiera andar, sobre hielo perfectamente sin \omterm{def:g10:inertia:inventory}{rozamiento}?
\end{exercise}

\begin{exercise}[$\star\star$]\label{exo:g11:forces-and-motion:10}
En una frenada de emergencia, ¿por qué un coche con ABS (con las ruedas
girando) suele parar antes que otro con las ruedas bloqueadas y
derrapando --- y por qué el coche que derrapa ya no se puede dirigir?
Cita la \cref{def:g11:forces-and-motion:friction}.
\end{exercise}

\begin{exercise}[$\star\star$]\label{exo:g11:forces-and-motion:11}
Un lanzador de martillo voltea la bola con su cable. ¿Qué \omterm{def:g10:inertia:force}{fuerza} hace
girar la bola? Esboza su camino tras el lanzamiento. ¿Por qué no
necesita cable la Luna, y en qué se convertiría su camino sin la
gravedad?
\end{exercise}

\begin{exercise}[$\star\star\star$]\label{exo:g11:forces-and-motion:12}
Una escoba es un palo uniforme de \qty{1.0}{kg} (con su propio \omterm{def:g11:forces-and-motion:com}{centro de
masas} a \qty{0.60}{m} del extremo superior) con un cabezal de
\qty{2.0}{kg} cuyo \omterm{def:g11:forces-and-motion:com}{centro de masas} está a \qty{1.20}{m} de ese extremo.
Tratando cada parte como una masa puntual, localiza el \omterm{def:g11:forces-and-motion:com}{centro de masas}
de la escoba y comprueba la regla de la razón inversa de la
\cref{def:g11:forces-and-motion:com}.
\end{exercise}

\begin{exercise}[$\star\star\star$]\label{exo:g11:forces-and-motion:13}
Un coche de masa $m$ y una furgoneta cargada de masa $2m$ frenan desde
la misma velocidad con la misma \omterm{def:g10:inertia:force}{fuerza} de frenado.
\begin{enumerate}
  \item Compara los tiempos que necesitan para pararse.
  \item Cada uno pierde velocidad de forma regular, así que ambos van a
        la misma velocidad \emph{media}: compara las distancias de
        frenado.
  \item Ahora el coche frena desde el \emph{doble} de velocidad, con la
        misma \omterm{def:g10:inertia:force}{fuerza}: compara el tiempo y la distancia con su primera
        frenada. ¿Moraleja para los límites de velocidad?
\end{enumerate}
\end{exercise}

\begin{exercise}[$\star\star\star$]\label{exo:g11:forces-and-motion:14}
La Luna orbita a $R \approx \qty{3.84e8}{m}$ en $T = \num{27.3}$ días.
\begin{enumerate}
  \item Calcula su velocidad orbital (la circunferencia entre el
        \omterm{def:g10:signals-and-waves:period}{periodo}).
  \item ¿Hacia dónde apunta su cambio de velocidad, y qué \omterm{def:g10:inertia:force}{fuerza} lo
        proporciona?
  \item Mejora, con la \cref{prop:g11:forces-and-motion:inward}, esta
        frase: ``la Luna no cae porque se mueve lo bastante deprisa''.
\end{enumerate}
\end{exercise}

\begin{exercise}[$\star\star\star$]\label{exo:g11:forces-and-motion:15}
Un huevo de masa \qty{50}{g} llega al suelo a \qty{4.0}{m/s}: sobre
baldosa se detiene en aproximadamente \qty{1}{ms} y sobre espuma gruesa,
en unos \qty{50}{ms}. Compara las \omterm{def:g10:inertia:force}{fuerzas} medias. Explica, con el mismo
razonamiento, por qué flexionas las rodillas al caer de un salto y de
qué está hecha la colchoneta de una gimnasta.
\end{exercise}

\section{Problema: el ensayo de choque}

\begin{problem}[{Problema de fin de semana --- diseñar la colisión
suave: puesto que un choque fija el cambio de velocidad, todo dispositivo
de seguridad de un coche es una máquina para estirar el tiempo}]
\label{pb:g11:forces-and-motion:1}
Un laboratorio de ensayos de choque lanza un coche a \qty{50}{km/h}
contra un muro rígido, con un maniquí adulto de \qty{70}{kg} y otro
infantil de \qty{20}{kg} a bordo; las cámaras cronometran cada parada.
Tu tarea: encontrar dónde vive la violencia de un choque y diseñarla
hasta hacerla desaparecer.

\medskip
\noindent\textbf{Parte I --- La medida de un choque.}
\begin{enumerate}
  \item Convierte \qty{50}{km/h} a \unit{m/s}.
  \item En cualquier parada frontal, el cambio de velocidad de los
        ocupantes tiene siempre el mismo módulo. ¿Cuál? ¿Por qué no
        puede alterarlo ningún dispositivo?
  \item Cita el \cref{thm:g11:forces-and-motion:secondlaw}: con
        $\Delta v$ y $m$ fijados, $F\,\Delta t$ está fijado. ¿Cuál es el
        único mando que le queda al diseñador?
  \item El muro detiene el coche en \qty{0.050}{s}; una zona de
        deformación estira eso hasta \qty{0.150}{s}. Compara las \omterm{def:g10:inertia:force}{fuerzas}
        medias.
  \item ¿Hacia dónde apuntan $\Delta \vect v$ y la \omterm{def:g11:forces-and-motion:netforce}{fuerza neta} sobre los
        maniquíes durante la parada?
\end{enumerate}

\medskip
\noindent\textbf{Parte II --- La zona de deformación.}
\begin{enumerate}[resume]
  \item ¿Por qué, en una frase, diseñan los ingenieros la parte
        delantera del coche para que se destruya?
  \item En un coche antiguo y rígido (parada en \qty{0.050}{s}), el
        maniquí está atado rígidamente al asiento. Compara su \omterm{def:g10:inertia:force}{fuerza} con
        la de la parada moderna de \qty{0.150}{s}.
  \item La velocidad del coche que se detiene cae de forma regular (de
        media, la mitad de la inicial). ¿Qué distancia recorre en
        \qty{0.150}{s}? Compárala con una zona de deformación real.
  \item ¿Por qué no una zona de deformación de \qty{10}{m}? ¿Y qué parte
        del coche \emph{no} debe deformarse?
  \item Resume la Parte II como un lema: ``no puedes elegir $\Delta v$,
        así que tienes que elegir \dots''.
\end{enumerate}

\medskip
\noindent\textbf{Parte III --- El cinturón, el airbag y el niño.}
\begin{enumerate}[resume]
  \item Un maniquí sin cinturón conserva sus \qty{14}{m/s}
        (\cref{ch:g10:inertia}) hasta que el parabrisas lo detiene en
        \qty{0.005}{s}; uno con cinturón se detiene con el coche en
        \qty{0.150}{s}. Compara las \omterm{def:g10:inertia:force}{fuerzas} medias.
  \item ¿Por qué es mejor un cinturón que se estira un poco que un arnés
        de acero inextensible?
  \item La cabeza se detendría sobre el volante en \qty{0.010}{s}; el
        airbag estira eso hasta \qty{0.080}{s}. Compara las \omterm{def:g10:inertia:force}{fuerzas}.
  \item El adulto y el niño sufren el mismo $\Delta v$ en el mismo
        $\Delta t$. Compara las \omterm{def:g11:forces-and-motion:netforce}{fuerzas netas} que deben aportar los dos
        cinturones.
  \item Un niño en el regazo: los brazos tienen que cambiar en
        \qty{0.150}{s} la velocidad de un bebé de \qty{10}{kg} en
        \qty{14}{m/s}, mientras que la gravedad cambia una velocidad en
        \qty{9.8}{m/s} cada segundo. ¿Cuántas veces más deprisa tienen
        que actuar los brazos --- sostener qué masa pesaría así?
        Concluye.
  \item Los niños sufren \omterm{def:g10:inertia:force}{fuerzas} menores --- entonces, ¿por qué las
        sillitas infantiles? Piensa en \emph{dónde} aplica su \omterm{def:g10:inertia:force}{fuerza} el
        cinturón, y en $\Delta t$.
\end{enumerate}

\medskip
\noindent\textbf{Parte IV --- La celda de seguridad y el veredicto.}
\begin{enumerate}[resume]
  \item El habitáculo es rígido mientras los dos extremos se deforman.
        ¿Por qué no debe deformarse el habitáculo, si la rigidez implica
        paradas violentas?
  \item Dos coches idénticos chocan de frente, cada uno a
        \qty{50}{km/h}. ¿Dónde se detiene cada uno? Compara cada choque
        con el ensayo contra el muro.
  \item Un coche pequeño choca de frente con un camión cargado. Las
        \omterm{def:g10:inertia:force}{fuerzas} que se ejercen mutuamente son iguales y opuestas, mutuas
        como toda interacción
        (\cref{ch:g11:fundamental-interactions}). ¿A quién le cambia más
        la velocidad?
  \item Enumera los dispositivos de seguridad de este problema y la
        única magnitud que estiran todos ellos.
  \item Final: enuncia en una frase la ley de diseño de la colisión
        suave usando únicamente \emph{masa}, \emph{cambio de velocidad},
        \emph{\omterm{def:g10:inertia:force}{fuerza}} y \emph{tiempo}.
\end{enumerate}
\end{problem}
